
\chapter*{Welcome from Colm Dunphy, the HDip Project Co-ordinator}

\vspace{-21pt}

%\def\WRAPWIDTH{0.225\textwidth}\def\WIDTH{0.2\textwidth}
%\begin{wrapfigure}[10]{r}{\WRAPWIDTH}
%\centering
%\vspace{-0.25cm} \fancyframe{\includegraphics[width=\WIDTH]
\def\WRAPWIDTH{4.5cm}\def\WIDTH{3.7cm}
\begin{wrapfigure}[12]{r}{\WRAPWIDTH}
\centering
\vspace{-0.25cm} \fancyframe{\includegraphics[clip,trim=40 40 40 40,width=\WIDTH]
{../common_images/orig/Colm_Dunphy_2026}}
\end{wrapfigure}

\vspace{0.4cm}
The multi award-winning Higher Diploma in Computer Science (online) was the first fully online programme from WIT, now SETU, and pioneered the innovative Agile Semester (multi-sprint) approach to delivery. The programme was designed to be delivered online, with an emphasis on student experience, engagement, and building a learner community that promotes peer learning and openness. Our graduates have a remote-first mentality, greater maturity from wider learning experiences, and strong capability. 

\vspace{0.4cm}
Students graduate as full-stack-oriented developers, primed with the foundations on which to build and exploit the new-gen AI, vibe-coded, spec-driven, and agentic AI software development world we now inhabit. Students on the programme complete their studies and a capstone project while on a six-month work placement/internship. This makes them available to be retained by your company (without returning). Students and graduates continue to be highly sought after. If your company is interested in providing a work placement/internship where the student will be mentored, please contact \href{mailto:joan.mangan@setu.ie}{joan.mangan@setu.ie}. 

\vspace{0.4cm}
This year’s projects span six broad categories: healthcare and life sciences, business and operational systems, AI and intelligent software, infrastructure and developer tooling, IoT and embedded systems, and consumer, civic, and creative apps. Across the portfolio, many projects are web-based, including backend, full-stack, and hybrid mobile/web applications, alongside work in systems, networking, security, MLOps, DevOps, cloud CI/CD, and IoT/physical computing. The collection also includes workplace projects built with platforms such as AWS, Azure, Linux, Kubernetes, and Red Hat OpenShift, as well as AI-driven solutions. Several projects contributed to open source, including FlightCtl. Many student projects integrate multiple APIs and deploy across different environments. Some projects have workplace details withheld under NDA.

\vspace{0.4cm}
For the HDip in Computer Science (online) section of this booklet, student thumbnails link to project videos, and student names link to the project page. Project pages summarise the project and include a link to a project landing page where additional links can be found to the GitHub repo, project videos on YouTube, the final report, the final presentation, and more. QR codes for each project are also provided.

\vspace{0.4cm}
In summary, we are often asked what our course is about and what you can do after completing it. Well, this showcase answers both questions through our students' hard work. Enjoy the diversity, innovation, and creation. From the entire team, we would like to thank our students for their work over the last few years, their tenacity and resilience in completing what is a flexible but intensive full-time programme delivered in an accelerated mode, part time, all the while working in a different discipline, preparing for a future in computing. From all the team, we wish you every success, which you will no doubt have in the future.

\vspace{0.4cm}
Regards,\\
Colm Dunphy 

\vspace{-0.3cm}

\clearpage
